\chapter{Abdominal Rigidity}

\section*{Definition}
Abdominal rigidity is involuntary, persistent tensing or hardening of the abdominal muscles. Unlike voluntary tightening or anxiety-related tensing, true rigidity remains when the person tries to relax. It is a physical sign of possible irritation of the abdominal lining (peritoneum), often from peritonitis, and requires urgent medical assessment.

\section*{When to See a Doctor}
True abdominal rigidity should be assessed urgently. Call emergency services or go to the emergency room immediately if rigidity occurs with:
\begin{itemize}
  \item Your abdomen feels hard or rigid and you have severe pain
  \item The pain is worse with movement, coughing, or being touched
  \item Nausea, vomiting, and inability to pass gas or have a bowel movement
  \item Fever, chills, or feeling severely unwell
  \item You are breathing rapidly or your heart is racing
  \item You feel dizzy, lightheaded, or faint (signs of internal bleeding or shock)
\end{itemize}

\section*{How It Develops}
When the lining of the abdominal cavity (the peritoneum) becomes irritated or inflamed, the body may respond by reflexively contracting the abdominal muscles. This protective response limits movement of the inflamed area. The abdomen may feel firm or board-like, and pain may worsen with movement, coughing, deep breathing, or examination. Voluntary muscle tensing caused by fear or anticipation can resemble guarding but should not be confused with true involuntary rigidity.

\section*{Causes}
\begin{itemize}
  \item \textbf{Peritonitis:} Inflammation or infection of the abdominal lining, often caused by leakage of intestinal contents, pus, blood, or other fluid
  \item \textbf{Appendicitis:} Inflammation of the appendix, especially if it has perforated
  \item \textbf{Perforated peptic ulcer:} A stomach or intestinal ulcer that has broken through the wall of the digestive tract
  \item \textbf{Diverticulitis:} Inflammation of diverticula, sometimes complicated by perforation or abscess
  \item \textbf{Cholecystitis or gallbladder perforation:} Severe inflammation or leakage from the gallbladder
  \item \textbf{Pancreatitis:} Severe inflammation of the pancreas, which can cause marked upper-abdominal pain and guarding
  \item \textbf{Ruptured ectopic pregnancy or ovarian cyst:} Internal bleeding or irritation in the pelvis
  \item \textbf{Abdominal trauma:} Injury to the abdomen that damages internal organs or causes bleeding
  \item \textbf{Surgical complications:} Infection, leakage, abscess, or bleeding after abdominal surgery
\end{itemize}

\section*{Related Symptoms}
\begin{itemize}
  \item Severe abdominal pain that worsens with movement
  \item Abdominal tenderness (pain when the abdomen is pressed)
  \item Rebound tenderness (pain when pressure is released suddenly)
  \item Nausea and vomiting
  \item Fever and chills
  \item Bloating or abdominal distension
  \item Inability to pass gas or have a bowel movement
  \item Rapid heart rate or breathing
\end{itemize}
\textbf{Important:} Abdominal rigidity is a sign of a serious medical condition, usually peritonitis or another source of peritoneal irritation. Do not eat or drink while arranging emergency assessment unless emergency clinicians advise otherwise.

\section*{Medical Evaluation}
\begin{itemize}
  \item A clinician will examine the abdomen gently for tenderness, involuntary guarding, rigidity, rebound, distension, hernias, and other signs of peritoneal irritation; a pelvic or rectal examination may be appropriate.
  \item Vital signs, mental status, urine output, and overall appearance are assessed for infection, dehydration, bleeding, or shock.
  \item Blood tests may include a complete blood count, electrolytes, kidney and liver tests, lipase, lactate, inflammatory markers, and blood cultures when sepsis is suspected. Urine and pregnancy testing are used when relevant.
  \item CT of the abdomen and pelvis is often the most informative imaging test for suspected perforation, abscess, obstruction, or other acute intra-abdominal disease. Chest or abdominal radiographs may show free air but can miss perforation; ultrasound is useful for selected biliary, pelvic, pregnancy, or fluid-related causes.
  \item Early surgical or specialist consultation is often needed. The cause and severity determine whether treatment requires surgery, interventional radiology, endoscopy, or medical management.
\end{itemize}

\section*{Treatment Options}
\begin{itemize}
  \item Source control is essential when there is perforation, abscess, infected tissue, or uncontrolled leakage. This may involve surgery, endoscopic treatment, or image-guided drainage.
  \item Intravenous (IV) broad-spectrum antibiotics are given promptly when bacterial peritonitis or sepsis is suspected, according to the likely source and local guidance.
  \item IV fluids are given to maintain hydration and blood pressure.
  \item Pain medication and medicines for nausea may be given during evaluation; appropriate analgesia does not need to be withheld while the cause is being investigated.
  \item Nasogastric tube (a tube through your nose into your stomach) may be placed to decompress the stomach and intestines.
  \item The specific treatment depends on the cause found during evaluation.
\end{itemize}

\section*{Self-Care and Home Management}
\begin{itemize}
  \item Do not eat or drink while arranging emergency assessment unless emergency clinicians advise otherwise; do not delay care because of concerns that pain medicine will mask the diagnosis.
  \item Do not apply heat or pressure to your abdomen or repeatedly test for rebound tenderness.
  \item Call emergency services or go to the nearest emergency room immediately, and do not drive yourself if you feel faint or severely unwell.
\end{itemize}

\section*{References}
\begin{thebibliography}{99}
\bibitem{abdominal_rigidity:ref1} Paterson-Brown S. ``The acute abdomen and surgical gastroenterology.'' In: \textit{Bailey \& Love\'s Short Practice of Surgery}, 28th ed. CRC Press, 2022.
\bibitem{abdominal_rigidity:ref2} Cartwright SL, Knudson MP. ``Evaluation of acute abdominal pain in adults.'' \textit{Am Fam Physician}. 2015;92(10):902-908.
\bibitem{abdominal_rigidity:ref3} Silen W. \textit{Cope\'s Early Diagnosis of the Acute Abdomen}, 22nd ed. Oxford University Press, 2010.
\bibitem{abdominal_rigidity:ref4} Debnath J, Sharma V, et al. ``Role of CT scan in the evaluation of acute abdomen.'' \textit{J Clin Diagn Res}. 2017;11(6):TE01-TE05.
\bibitem{abdominal_rigidity:ref5} World Society of Emergency Surgery. ``WSES guidelines on the management of acute peritonitis.'' \textit{World J Emerg Surg}. 2020;15:28.
\bibitem{abdominal_rigidity:ref6} MedlinePlus. Abdominal Rigidity; Peritonitis; Gastrointestinal Perforation. U.S. National Library of Medicine. Accessed 2026.
\end{thebibliography}
